\documentclass[12pt]{article}
\usepackage[utf8]{inputenc}
\usepackage[T1]{fontenc}
\usepackage{hyperref}
\title{Waray-Waray Language Resources for K-12 Curriculum Development}
\author{Compiled by Claude Sonnet 5 (claude-sonnet-5) for Andr\'as Kornai}
\date{}
\begin{document}
\maketitle

\section{Introduction}
Waray-Waray (ISO 639-3 \texttt{war}, Glottocode wara1300) is an Austronesian language of the Eastern Visayas in the Philippines, with about 3.4 million speakers. It is used as a language of instruction under the Philippine mother-tongue-based multilingual education (MTB-MLE) policy. This report lists only resources that appeared in web search results. It is a short survey, not an exhaustive one.

\section{Primary Educational Resources}
\begin{itemize}
\item DepEd LRMDS, MTB-MLE Grade 1 Waray, Week 10 Learner's Material: \url{https://lrmds.deped.gov.ph/detail/4948}
\item DepEd LRMDS, MTB-MLE Grade 1 Waray, Week 12 Teacher's Guide: \url{https://lrmds.deped.gov.ph/detail/4993}
\end{itemize}
The search found only Grade 1 materials. I did not find the higher grades.

\section{Linguistic and Reference Materials}
\begin{itemize}
\item \emph{Diksiyonaryo Caton: Waray}, a Vibal trilingual dictionary series aligned with MTB-MLE: \url{https://www.philippinebooks.com/products/diksiyonaryo-caton-waray}
\item Orejola, \emph{Learning Waray-Waray}, a Waray-English and English-Waray dictionary, Vol.~2 (commercial listing; availability is uncertain): \url{https://origin-www.alibris.com/Learning-Waray-Waray-English-and-English-Waray-Dictionary-Vol-2-Dr-Orejola/book/52954064?qsort=p}
\item \emph{Bisayan grammar and notes on Bisayan rhetoric and poetics} (public-domain historical grammar): \url{https://gutenberg.org/ebooks/69603}
\item Swarthmore LING073 wiki page on Waray: \url{https://wikis.swarthmore.edu/ling073/Waray}
\item A Scribd upload of a short Waray dictionary. Its provenance is unverified, so treat it with caution: \url{https://www.scribd.com/doc/269809110/Waray}
\end{itemize}

\section{Religious and Community Resources}
None were found in this search.

\section{Digital and Audio Resources}
\begin{itemize}
\item IJERE article on MTB-MLE initiatives in Region 8, which describes a Waray corpus and frequency lists: \url{https://ijere.iaescore.com/index.php/IJERE/article/view/4462}
\item University of the Philippines journal article, found in search; I did not review its contents: \url{https://journals.upd.edu.ph/index.php/archive/article/download/9334/8248}
\end{itemize}

\section{Curriculum Development Context}
DepEd's MTB-MLE programme provides Waray learning materials, at least for the early grades. A curriculum developer should check the DepEd LRMDS portal for the other grades and for materials in subjects other than the language itself.

\end{document}